\chapter{Introdução}
\label{ch:introducao}

O mercado de ações desempenha um papel crucial na economia, sendo o comportamento dos preços das ações um indicador essencial para avaliar o momento econômico de empresas e setores. Essas variações refletem o sentimento dos investidores, influenciando diretamente o retorno sobre o capital investido (\cite{Smales2017}).

Entretanto, a característica marcante do mercado de ações é sua volatilidade, evidenciada por mudanças frequentes no cenário. A pandemia de COVID-19 (2020-2022), um evento imprevisível, representou um dos períodos mais voláteis na história recente do mercado de ações. A desconfiança dos investidores em relação à duração da pandemia e à capacidade das empresas de manterem a lucratividade diante da redução de receitas resultou em uma queda generalizada nos preços das ações em todo o mundo (\cite{Wang2021}).

Embora mudanças de cenário muitas vezes sejam imprevisíveis, certas variáveis, como o histórico de preços de uma ação ao longo do tempo, podem fornecer \textit{insights}, ou seja, novas perspectivas valiosas para antecipar movimentos futuros.

Nos últimos anos, observou-se no Brasil um crescente interesse por investimentos, incluindo o mercado de ações. Como muitas pessoas estão lidando pela primeira vez com esse universo, é natural que surjam dúvidas, muitas vezes relacionadas à volatilidade do mercado.

A área de \textit{Machine Learning}, que significa aprendizado de máquina em português, oferece ferramentas poderosas para analisar grandes conjuntos de dados e identificar padrões. Em campos como a saúde, já se explorou o uso de \textit{Machine Learning} para prever cenários, como doenças cardíacas e Alzheimer (\cite{Alotaibi2019, Antor2019}). A aplicação desses conceitos pode ser estendida ao mercado financeiro, onde técnicas de \textit{Machine Learning} têm sido empregadas para prever valores de ações com base em seus históricos de preços (\cite{Vijh2020}).

Este trabalho propõe utilizar técnicas de \textit{Machine Learning} para prever os valores de ativos no mercado financeiro. Reconhecemos a complexidade do ambiente, com inúmeras variáveis influenciando os preços das ações, como o cenário político, a imagem pública da empresa e crises econômicas. No entanto, acreditamos que fatores como histórico de preços, volume de negociações e indicadores de análise técnica contribuem para uma maior previsibilidade no curto prazo. A avaliação da eficácia dessa abordagem será realizada comparando os resultados previstos pela rede neural com os valores reais das ações, considerando também métricas de desempenho (\cite{Patel2015}).

\section{Objetivo}
\label{sec:objetivo}

Este estudo tem como objetivo principal a comparação entre modelos de previsão de valores das ações no mercado financeiro, adotando uma abordagem fundamentada na análise de dados históricos e variáveis macroeconômicas, fazendo uso de técnicas de \textit{Machine Learning}. Buscamos não apenas a acurácia nas previsões, mas também a compreensibilidade das saídas do modelo com a intenção de de fornecer à investidores previsões confiáveis e interpretações claras, capacitando-os a tomar decisões de investimento embasadas em um ambiente dinâmico e desafiador do mercado financeiro.

\subsection{Objetivos Específicos}
\label{subsec:objespecificos}

A fim de alcançar o objetivo geral, os seguintes objetivos específicos foram estabelecidos:

\begin{itemize}
    \item Elaboração de um método computacional que contribua para a tomada de decisões mais informadas no campo dos investimentos financeiros e para a compreensão dos padrões subjacentes aos mercados. Este objetivo será alcançado através do desenvolvimento e implementação de um modelo de previsão que seja capaz de antecipar os movimentos dos preços dos ativos com base em dados históricos.
    \item Comparar modelos de diferentes topologias para entender qual das diferentes abordagens oferece maior capacidade de generalização para o domínio do problema proposto.
    \item Realizar uma análise empírica para determinar a saída mais apropriada do modelo, visando fornecer previsões compreensíveis para os investidores. O objetivo é identificar o formato de saída que permita uma compreensão intuitiva das previsões, auxiliando efetivamente nas decisões de investimento.
\end{itemize}


\section{Organização do Trabalho}
\label{sec:organizacao}

O Capítulo 1, denominado "Introdução", contextualiza o escopo desta pesquisa, destacando o problema em foco, os objetivos da investigação e a abordagem metodológica adotada para alcançar esses objetivos.

No Capítulo 2, "Referencial Teórico", são apresentados os fundamentos teóricos, explorando conceitos de Inteligência Artificial, Machine Learning, Redes Neurais e Deep Learning. Este capítulo serve como base para a compreensão dos conceitos fundamentais que sustentam a implementação proposta.

O Capítulo 3, intitulado "Trabalhos Relacionados", examina estudos e pesquisas prévias relacionadas, identificando contribuições relevantes no contexto da previsão de valores de ações no mercado financeiro. Essa análise serve como referência para posicionar a proposta de maneira contextualizada no cenário acadêmico.

O Capitulo 4, "Metodologia", explica os processos de coleta e pré-processamento dos dados e também o treinamento desses dados. Este capítulo explica os critérios usados para a coleta dos dados e a importância desta etapa para o bom treinamento dos modelos propostos. 

No Capítulo 5, "Modelos Desenvolvidos", detalhamos os modelos utilizados para o treinamento da rede neural artificial. Este capítulo expõe escolhas técnicas e implementações específicas que embasam as decisões tomadas durante o processo de desenvolvimento.

O Capítulo 6, "Resultados Obtidos", avalia os desempenhos apresentados pelos diferentes modelos treinados, avaliação de qual modelo é mais aplicável ao mundo real. Considerando a previsão de valores de ações no mercado financeiro, delineando as contribuições e avanços alcançados.